\section{Discussion}
\label{sec:discussion}

\paragraph{What LNC is.} LNC is a reference-free metric that scores all four SummEval dimensions from two small open models at the cost of BERTScore. In domain it is more accurate on the shortcut-controlled axis than every other reference-free metric we could run, including a local LLM judge that costs 82 times more, and its consistency scorer transfers to FRANK, where it matches AlignScore. Its novelty is empirical rather than conceptual: every signal it uses has prior art, and the combination is a calibrated ridge regression. What is new is the choice of signals by their behaviour on the partialled axis, the per-dimension calibration, and the evidence about where such a calibration does and does not transfer.

\paragraph{What LNC is not.} LNC is not a replacement for UniEval when a reference exists: UniEval is better on coherence and fluency, uses a reference for relevance, and costs only 1.6 times more. It is not a general relevance metric: its relevance scorer does not transfer to content-coverage annotations. And it is not independent of SummEval: part of its calibration is specific to the systems it was fitted on, and a ridge over existing reference-based metrics reaches the same in-domain score. We therefore present it as a cheap reference-free metric for consistency, fluency and coherence of news summaries, with relevance as its weak dimension.

\paragraph{Recommendations for metric evaluation.} Three practices would have changed our own conclusions had we adopted them earlier. First, report correlations with extractiveness partialled out next to the raw ones, and include a lead-window control and the copy rate itself in every comparison; on SummEval they beat most metrics on the raw axis. Second, evaluate any metric that is calibrated on human ratings on held-out articles and, if possible, held-out systems and corpora; the dataset-order split of SummEval was favourable to LNC, and unseen systems were harder for it. Third, fix the transfer protocol before running it, and report the verdict even when it is negative.

\subsection{Limitations}

\paragraph{Baselines below their published strength.} Our G-Eval judge is a quantised 7B model with three samples, not GPT-4 with twenty, and GPTScore uses GPT-2 small. Both are therefore weaker than in their original papers, and the comparisons with them characterise cheap local configurations, not the methods in general. Our UniEval omits the sentence-level aggregation of the original for consistency, which may explain part of its weakness on that dimension. AlignScore uses the official weights and is the fairest of the strong baselines.

\paragraph{One calibration corpus and one language.} LNC is calibrated on 800 English news summaries rated by three experts each. Its features are language-specific in parts (the perturbation rules and the entity heuristic), and its transfer was tested only on English news and XSum.

\paragraph{The partial axis is a single control.} Partialling out the bigram copy rate removes one linear, rank-level confound. It does not remove nonlinear dependence on copying or other confounds such as length, and the copy rate is itself correlated with quality: a summary that copies accurately is often a good summary. The partial axis is a diagnostic of which part of a metric's agreement is not explained by copying, not a definition of quality.

\paragraph{Cost on one machine.} Wall-clock cost depends on the hardware, batching and runtime. We measured every metric on the same machine with the models preloaded, which makes the comparison internally fair, but absolute numbers will differ elsewhere. The ratios we emphasise span one to four orders of magnitude, larger than the variation between typical GPU setups.

\paragraph{No downstream validation.} We did not test LNC as a reward for training, as a selection criterion between systems in deployment, or for the interpretability of its feature weights by human readers. Its system-level correlations (Appendix~\ref{app:tables}) are high but rest on 16 systems.

\subsection{Threats to validity}

The SummEval development half was used for every design decision, but the signals that LNC combines were found during a search in which we looked at full-set correlations of earlier metrics, so the held-out half is not entirely untouched by our choices. The random-split and unseen-system analyses, which refit the calibration from scratch, and the frozen transfer test bound the effect of this. The transfer corpora were converted to SummEval's layout by our own script; we release it, and it rebuilds the files byte for byte from their public sources. All reported numbers are generated by scripts from stored per-sample scores, and a second run of the whole pipeline on the same machine reproduced the scores of LNC and of every deterministic metric exactly.

\section{Conclusion}
\label{sec:conclusion}

LNC shows that two small open models are enough for a reference-free metric that covers all four dimensions of summary quality: a calibrated combination of language-model likelihood, perturbation contrasts and sentence-level entailment is cheaper than BERTScore and more accurate than a local LLM judge, and it keeps most of its accuracy when the extractive shortcut of the benchmark is controlled for. A pre-registered test shows that its consistency judgement transfers to new systems and annotation schemes while its relevance judgement does not, and controls show where its in-domain accuracy comes from. We hope that the evaluation protocol, with shortcut-controlled correlations, same-machine costs, frozen transfer and calibrated controls, proves as useful as the metric.

\section*{Data and code availability}
All code, the four corpora in a common format, the per-sample scores of every metric, the frozen LNC calibration and the scripts that render every table in this paper are available at \url{https://github.com/mikolajsemeniuk/llmbench}. The command \texttt{make clean data reproduce} rebuilds the corpora from their public sources and regenerates every result. SummEval, FRANK and RoSE are distributed by their authors under their own licences.

\section*{CRediT authorship contribution statement}
\textbf{Mikolaj Semeniuk:} Conceptualization, Methodology, Software, Validation, Formal analysis, Investigation, Data curation, Writing -- original draft. \textbf{Agnieszka Jastrzebska:} Conceptualization, Methodology, Validation, Supervision, Writing -- review \& editing.

\section*{Declaration of competing interest}
The authors declare that they have no known competing financial interests or personal relationships that could have appeared to influence the work reported in this paper.

\section*{Funding}
This research did not receive any specific grant from funding agencies in the public, commercial, or not-for-profit sectors.

\section*{Declaration of generative AI and AI-assisted technologies in the writing process}
During the preparation of this work the authors used Claude (Anthropic) in order to assist with writing analysis code and with drafting and editing the text. After using this tool, the authors reviewed and edited the content as needed and take full responsibility for the content of the publication.
